\documentclass[11pt]{amsart}
\usepackage[margin=1.1in]{geometry}
\usepackage[T1]{fontenc}
\usepackage{amsmath,amssymb,amsthm,mathtools,enumitem}
\usepackage[hidelinks]{hyperref}

\newtheorem{theorem}{Theorem}[section]
\newtheorem{proposition}[theorem]{Proposition}
\newtheorem{lemma}[theorem]{Lemma}
\newtheorem{corollary}[theorem]{Corollary}
\theoremstyle{definition}
\newtheorem{definition}[theorem]{Definition}
\theoremstyle{remark}
\newtheorem{remark}[theorem]{Remark}
\numberwithin{equation}{section}

\newcommand{\Th}{\text{\th}}
\newcommand{\cL}{\mathcal{L}}
\newcommand{\R}{\mathbb{R}}
\newcommand{\C}{\mathbb{C}}
\newcommand{\Z}{\mathbb{Z}}
\renewcommand{\Im}{\operatorname{Im}}
\renewcommand{\Re}{\operatorname{Re}}

\title[Gravitational mode stability of Kerr--de Sitter]{Mode stability for gravitational perturbations of subextremal Kerr--de Sitter black holes}
\author{Jaideep Sai Padhi}
\address{Purdue University, West Lafayette, IN, USA}
\email{jpadhi@purdue.edu}
\date{November 8, 2026}

\begin{document}

\begin{abstract}
We prove mode stability for the Teukolsky equations of spin $s=\pm2$ on Kerr--de Sitter black hole spacetimes in the
full subextremal range: there are no nontrivial outgoing mode solutions at frequencies $\sigma$ with
$\Im\sigma\ge0$. We then show that every generalized mode solution of the ungauged linearized Einstein equation
$D_g\mathrm{Ric}(h)=\Lambda h$ with $\Im\sigma\ge0$ is pure gauge if $\sigma\ne0$, and the sum of a linearized
Kerr--de Sitter metric and a pure gauge term if $\sigma=0$. This is the mode stability hypothesis under which
Hintz, Petersen and Vasy proved the nonlinear stability of Kerr--de Sitter black holes in the full subextremal range.
For real frequencies, the Teukolsky part follows Hintz's two-current method for scalar waves, applied to a dual
radial equation obtained from an Euler integral transform. The upper half plane is reached by a continuity argument
in which high-frequency control is borrowed from the gauge-fixed Einstein operator through metric reconstruction.
The passage from the Teukolsky scalars to the metric uses a linearized Goldberg--Sachs argument and the two Killing
fields built from the Killing spinor. Several steps are computer-assisted, using exact arithmetic.
\end{abstract}

\maketitle
\tableofcontents

\section{Introduction}\label{sec:intro}

For $\Lambda>0$, $M>0$ and $a\in\R$, the Kerr--de Sitter (KdS) metric in Boyer--Lindquist coordinates is
\begin{equation}\label{eq:kds}
g=-\frac{\Delta_r}{\Xi^2\varrho^2}\big(dt-a\sin^2\theta\,d\varphi\big)^2+\varrho^2\Big(\frac{dr^2}{\Delta_r}+\frac{d\theta^2}{\Delta_\theta}\Big)
+\frac{\Delta_\theta\sin^2\theta}{\Xi^2\varrho^2}\big((r^2+a^2)\,d\varphi-a\,dt\big)^2,
\end{equation}
where
\[
\Delta_r=(r^2+a^2)\Big(1-\frac{\Lambda r^2}{3}\Big)-2Mr,\quad \varrho^2=r^2+a^2\cos^2\theta,\quad
\Delta_\theta=1+\alpha\cos^2\theta,\quad \alpha=\frac{\Lambda a^2}{3},\quad \Xi=1+\alpha.
\]
It solves $\mathrm{Ric}(g)=\Lambda g$. The parameters are \emph{subextremal} if $\Delta_r$ has four distinct real
roots, which we label
\[
r_3<0\le r_0<r_1<r_2,
\]
where $r_0=r_C$ is the Cauchy horizon, $r_1=r_e$ the event horizon and $r_2=r_c$ the cosmological horizon;
then $r_0+r_1+r_2+r_3=0$. The domain of outer communications is $r_1<r<r_2$. We use a time function $t_*$ and an
azimuthal angle $\varphi_*$ that are smooth across the future event and cosmological horizons, as in
\cite{HPV25,Hin26}; a \emph{mode} at frequency $\sigma\in\C$ is a function, or a tensor, of the form
$e^{-i\sigma t_*}v$ with $v$ independent of $t_*$ and smooth on $[r_1-\delta,r_2+\delta]\times S^2$. In the language of
\cite{Hin26}, such modes are \emph{outgoing}.

The nonlinear stability of the KdS family was proved by Hintz and Vasy \cite{HV18} for very slowly rotating black
holes. Hintz, Petersen and Vasy \cite{HPV25} extended this to the full subextremal range, conditional on mode
stability for the ungauged linearized Einstein equation \cite[Assumption~1.2]{HPV25}. Hintz \cite{Hin26} recently
proved mode stability for scalar waves in the full subextremal range, and left the Teukolsky equations of spin
$s=\pm2$ open. Our main results settle both questions.

\begin{theorem}[Teukolsky mode stability]\label{thm:A}
On every subextremal KdS spacetime, the Teukolsky equations of spin $s=+2$ and $s=-2$ have no nontrivial outgoing
mode solutions at any frequency $\sigma\in\C$ with $\Im\sigma\ge0$, including $\sigma=0$. The same holds for
generalized modes $\sum_{j<k}t_*^je^{-i\sigma t_*}v_j$.
\end{theorem}

\begin{theorem}[Ungauged linearized Einstein mode stability]\label{thm:B}
Let $(\Lambda,M,a)$ be subextremal with $a\neq0$, and let $T=\partial_{t_*}+\frac{a}{r_2^2+a^2}\partial_{\varphi_*}$.
Let $h$ be a smooth symmetric 2-tensor with $D_g\mathrm{Ric}(h)=\Lambda h$.
\begin{enumerate}[label=(\arabic*)]
\item If $(\cL_T+i\sigma)h=0$ with $\sigma\ne0$ and $\Im\sigma\ge0$, then $h=\cL_Xg$ for a smooth vector field $X$.
\item If $(\cL_T)^kh=0$ for some $k\ge1$, then $h=c_M\partial_Mg+c_a\partial_ag+\cL_Xg$ with $c_M,c_a\in\C$ and $X$
smooth.
\end{enumerate}
\end{theorem}

For $a=0$, Theorem~\ref{thm:B} is \cite[\S7]{HV18}. Theorem~\ref{thm:B} is \cite[Assumption~1.2]{HPV25} (with $\cL_Xg=2\delta^*X^\flat$),
so \cite[Theorem~1.4]{HPV25} now applies unconditionally:

\begin{corollary}
Subextremal Kerr--de Sitter black holes are nonlinearly stable as solutions of the Einstein vacuum equations with
positive cosmological constant, in the sense of \cite[Theorem~1.4]{HPV25}.
\end{corollary}

\subsection{Strategy}
The proof of Theorem~\ref{thm:A} has three parts.

\emph{Real frequencies} (\S\S\ref{sec:radial}--\ref{sec:event}). After separation of variables, outside the
superradiant range a Wronskian identity suffices. Inside it, we do not work with the radial Teukolsky equation
itself: we apply an Euler (Riemann--Liouville) integral transform to its Heun form. The transform is injective, and its
image solves a dual Heun equation that is \emph{real}. Its potential has the very simple form
\[
L(z)=\frac{A}{z}+B-3z,\qquad A>0,
\]
in a variable $z$ in which the event and cosmological horizons sit at $z=1$ and $z=t$. The dual equation has exactly
the structure of the scalar radial equation in \cite{Hin26}, so Hintz's two-current method applies. It reduces mode
stability to an \emph{event lemma}: an ordinary differential inequality in the four parameters $(t,z_0,A,B)$. A lower
bound on the spin-weighted spheroidal eigenvalue gives the a priori bound $B\ge5.5t-1.25$. We prove the event lemma
in three regimes. In the hardest one we follow \cite[\S4.3.3]{Hin26}: we first saturate one of the two conditions and
then solve the other with equality. The remaining claims reduce to sign statements about explicit polynomials, which
we certify by exact Bernstein-basis branch and bound.

\emph{Zero frequency} (\S\ref{sec:zero}). This is an energy identity combined with a Teukolsky--Starobinsky map.

\emph{The upper half plane} (\S\ref{sec:uhp}). We argue by continuity along a path of subextremal parameters ending
at Schwarzschild--de Sitter (SdS), as in \cite[\S6]{Hin26}. The new difficulty compared with scalar waves is that no
high-frequency estimate for the Teukolsky operator in the full subextremal range is available. We circumvent this.
\begin{itemize}
\item Using a Teukolsky mode as an outgoing-radiation-gauge Hertz potential \cite{Wal78} produces a linearized
Einstein mode with nonvanishing $\Psi_4$.
\item Gauge fixing turns it into a quasinormal mode of the gauge-fixed operator of \cite{HPV25}.
\item The high-energy estimates of \cite{HPV25} then confine such modes to a compact set.
\end{itemize}
Nonvanishing of $\Psi_4$ uses an elementary observation: away from an explicit discrete set of frequencies, nothing
annihilated by $(\Th')^4$ can be regular at \emph{both} horizons. This uses the cosmological horizon in an essential
way.

\emph{Theorem~\ref{thm:B}} (\S\ref{sec:B}). Once the Teukolsky scalars vanish, a linearized Goldberg--Sachs argument
shows that the perturbed spacetime is algebraically special to first order. Two Killing fields can then be built
covariantly from the perturbed Weyl tensor: the field $\xi$ obtained from the Killing spinor, and Carter's field
$\zeta$. Both remain Killing to first order. Moreover, $\Im\delta\xi$ is a Killing field of the background, so it
vanishes for every mode with $(\sigma,m)\ne(0,0)$. This makes $h$ pure gauge, with an explicit gauge vector. The
stationary axisymmetric sector reduces to Carter's separable ansatz, on which the field equations are linear.

\subsection*{Computer assistance}
Some steps rely on exact symbolic computation or on certificates in exact rational arithmetic. They are listed in
Appendix~\ref{app:cert}. No step relies on floating-point evidence.

\subsection*{Acknowledgments}
I thank the authors of \cite{Hin26,HPV25} for making their work publicly available.

\section{Preliminaries}\label{sec:prelim}

\subsection{Parameters}
By Vieta's formulas,
\begin{equation}\label{eq:vieta}
\sum_j r_j=0,\qquad \sum_{i<j}r_ir_j=-\frac{3(1-\alpha)}{\Lambda},\qquad r_0r_1r_2r_3=-\frac{3a^2}{\Lambda}.
\end{equation}
We write $\rho_j=r_j^2+a^2$, $\Omega_j=a/\rho_j$, $D_j=|\Delta_r'(r_j)|$ and $P_j=\rho_j/D_j$. The signs of
$\Delta_r'(r_j)$ are $(-,+,-,+)$ for $j=0,1,2,3$. Following \cite[\S6, footnote~5]{Hin26}, we use the roots as
parameters.

\begin{lemma}[Rotation bound]\label{lem:alpha}
On every subextremal KdS spacetime, $\alpha<7-4\sqrt3$.
\end{lemma}
\begin{proof}
By scaling we may assume $\Lambda=3$. Then $-\Delta_r=r^4-(1-\alpha)r^2+2Mr-\alpha$, so $p_2:=\sum r_j^2=2(1-\alpha)$
and $\mathcal A:=-r_0r_1r_2r_3=\alpha$. Scaling a configuration of roots by $\tau>0$ multiplies $\mathcal A$ by $\tau^4$ and
$p_2$ by $\tau^2$. Hence $\alpha$ is the positive root of $\alpha=1-\frac12\sqrt{\alpha/k}$, where $k:=\mathcal A/p_2^2$
is scale invariant, and this root is increasing in $k$. Put $S=-r_3=r_0+r_1+r_2$. By the AM--GM inequality,
$r_0r_1r_2\le S^3/27$, and by the QM--AM inequality, $r_0^2+r_1^2+r_2^2\ge S^2/3$. Thus
$k\le(S^4/27)/(16S^4/9)=1/48$, with equality only if $r_0=r_1=r_2$, which subextremality excludes. At $k=1/48$ the
root is $\alpha=12(3^{-1/2}-\frac12)^2=7-4\sqrt3$.
\end{proof}

\subsection{The Teukolsky equations and separation of variables}
For $\sigma$ real, the Teukolsky equation of spin $s$ separates as $e^{-i\sigma t+im\varphi}R(r)S(\theta)$
\cite{Kha83}. The angular operator is then symmetric with compact resolvent. As in \cite[\S3]{Hin26}, a mode solution
can therefore be decomposed into separated ones, each with an eigenvalue $\lambda\in\R$. We only separate variables
for real $\sigma$. Put
\[
K(r)=\Xi\big((r^2+a^2)\sigma-am\big),\qquad \kappa_j=\frac{iK(r_j)}{\Delta_r'(r_j)}.
\]
The residue theorem applied to $\rho/\Delta_r$, $1/\Delta_r$ and $r^k/\Delta_r$ gives
\begin{equation}\label{eq:residues}
\sum_j\kappa_j=0,\qquad\sum_jr_j\kappa_j=-i\Xi\sigma,\qquad\sum_jr_j^2\kappa_j=0,
\end{equation}
together with $P_1+P_3=P_0+P_2$ and $D_1^{-1}+D_3^{-1}=D_0^{-1}+D_2^{-1}$.

\begin{lemma}[Boundary conditions]\label{lem:bc}
A separated mode is outgoing if and only if
\[
R\in(r-r_1)^{-\kappa_1-s}\,C^\infty([r_1,r_2))\cap(r_2-r)^{\kappa_2}\,C^\infty((r_1,r_2]).
\]
\end{lemma}
\begin{proof}
This follows by writing $t=t_*+T(r)$ and $\varphi=\varphi_*+\Phi(r)$, with the singular parts of $T$ and $\Phi$ at
$r_1$ and $r_2$ as in \cite[(3.4)]{Hin26}. The additional factor $\Delta_r^{-s}$ at $r_1$ comes from the boost
relating the Kinnersley tetrad to a tetrad that is regular across the event horizon. The Kinnersley tetrad is
already regular across the future cosmological horizon in outgoing coordinates.
\end{proof}

\subsection{Heun form}
The M\"obius map $z$ with $z(r_0,r_1,r_2,r_3)=(0,1,t,\infty)$, where
\[
t=\frac{(r_1-r_3)(r_2-r_0)}{(r_1-r_0)(r_2-r_3)}>1,
\]
turns the radial Teukolsky equation into a Heun equation with singular points $0,1,t,\infty$. We write
$\mu(z)=z(z-1)(t-z)$. At $z_j$ ($j=0,1,t$) the exponents are those of Lemma~\ref{lem:bc}, and at $\infty$ they are
$\{1+2s,\,1+s-2\kappa_3\}$. Two constants will appear repeatedly:
\begin{equation}\label{eq:C0}
C_0=\frac{3}{\Lambda(r_1-r_0)(r_2-r_3)},\qquad c_1=1+t+3(1-\alpha)C_0.
\end{equation}

\section{Real frequencies: the radial problem}\label{sec:radial}

Throughout this section $\sigma\in\R\setminus\{0\}$. We fix a separated outgoing mode with eigenvalue $\lambda$. By
the symmetry $(\sigma,m,a)\mapsto(-\sigma,-m,a)$ we may assume $m\ge0$.

\subsection{The superradiant window}
Set $e_1=2(\kappa_1+\kappa_3)$ and $e_t=2(\kappa_2+\kappa_3)$. For real $\sigma$ the $\kappa_j$ are purely imaginary.

\begin{proposition}[The window]\label{prop:window}
Let $m>0$ and $w=\sigma/m$. Then $\Im e_1=2\Xi mf_1(w)$ and $\Im e_t=2\Xi mf_2(w)$, where
\[
f_1(w)=P_1(w-\Omega_1)+P_3(w-\Omega_3),\qquad f_2(w)=-P_2(w-\Omega_2)+P_3(w-\Omega_3).
\]
The function $f_1$ is increasing and $f_2$ is decreasing, and their zeros satisfy $0<w_2<w_1$. Consequently
$\Im e_1\cdot\Im e_t>0$ exactly on the open interval $(w_2,w_1)$, the \emph{window}.
\end{proposition}

\begin{proof}
The formulas follow from the definition of $\kappa_j$ and the signs of $\Delta_r'(r_j)$.
\begin{itemize}
\item \emph{Monotonicity.} The slope of $f_1$ is $P_1+P_3>0$. By the residue identities, the slope of $f_2$ is
$P_3-P_2=P_0-P_1$. Moreover $D_0/D_1=f(r_0)/f(r_1)$ for the concave parabola $f(x)=(x-r_3)(r_2-x)$, whose vertex
$-(r_0+r_1)/2$ lies to the left of $r_0$. Hence $D_0>D_1$, and with $\rho_1>\rho_0$ this gives $P_1>P_0$.
\item \emph{Positivity of $w_2$.} The same comparison gives $w_2=a(D_1^{-1}-D_0^{-1})/(P_1-P_0)>0$.
\item \emph{Ordering $w_2<w_1$.} Cross-multiplying, the terms involving $a^2$ cancel, and $D_0D_1D_3$ times the
difference $w_1-w_2$ factors as $2(r_0-r_1)(r_0-r_3)(r_1-r_3)(r_0r_3-r_1r_2)$. All four factors are negative,
since $r_0r_3\le0<r_1r_2$, so the product is positive. \qedhere
\end{itemize}
\end{proof}

For every real $\sigma\ne0$, Proposition~\ref{prop:dual} below puts the dual equation in the real form
\eqref{eq:dual}, and $u$ is outgoing at both ends. Let $B_1,B_t$ be the boundary coefficients of $u$ at $z=1,t$ in
the sense of \cite[(3.5)]{Hin26}. Conservation of $\mu\Im(u'\bar u)$ gives, as in \cite[\S3]{Hin26},
\begin{equation}\label{eq:wronskian}
\Im e_1\,|B_1|^2=t\,\Im e_t\,|B_t|^2.
\end{equation}

\begin{corollary}\label{cor:outside}
There are no outgoing modes with $\sigma\in\R\setminus\{0\}$ outside the open window, nor any with $m=0$.
\end{corollary}
\begin{proof}
For $m=0$, $\Im e_1=2\Xi\sigma(P_1+P_3)$ and $\Im e_t=2\Xi\sigma(P_3-P_2)=-2\Xi\sigma(P_1-P_0)$ have opposite signs.
Off the closed window the same holds by Proposition~\ref{prop:window}, so \eqref{eq:wronskian} forces
$B_1=B_t=0$. At the endpoints exactly one of $\Im e_1,\Im e_t$ vanishes, since $w_1\ne w_2$, and
\eqref{eq:wronskian} gives $B_t=0$ or $B_1=0$ at the other end. In each case the indicial roots at that end differ by
the non-integer $e_t$ or $e_1$, so $u=0$ as after \cite[(3.6)]{Hin26}. Then $\tilde y=0$, which contradicts
Proposition~\ref{prop:transform}.
\end{proof}

\subsection{The dual equation}\label{subsec:dual}
\begin{proposition}[Euler transform]\label{prop:transform}
Let $\sigma\in\R\setminus\{0\}$, and let $y\ne0$ be the Heun-form radial function of an outgoing mode, normalized so
that it is analytic at $z=t$ and lies in $(z-1)^{\rho_1}C^\omega$ near $z=1$, where $\rho_1=-2\kappa_1-s$. Put
$\nu=s-1-2\kappa_3$. Then there is a solution $\tilde y\ne0$ of a Heun equation $\widetilde{\mathcal H}\tilde y=0$,
obtained from $y$ by a Riemann--Liouville integral of order $\nu+1$ based at $z=1$, with the following properties.
\begin{enumerate}[label=(\roman*)]
\item The exponents of $\widetilde{\mathcal H}$ at $z_j\in\{0,1,t\}$ are $0$ and $-2(\kappa_j+\kappa_3)$.
\item $\tilde y$ is analytic at $z=t$, and near $z=1$ it lies in $(z-1)^{-e_1}C^\omega$, where $e_1=2(\kappa_1+\kappa_3)$.
\end{enumerate}
\end{proposition}

\begin{proof}
That the Euler transform with kernel $(z-w)^{\nu}$ maps solutions of a Heun equation to solutions of a Heun
equation, when $\nu+2$ is an exponent at $\infty$, is classical \cite{KS96,Tak03}; the exponents in (i) follow from
those at $\infty$ listed in \S\ref{sec:prelim}. Write $I^{\nu+1}$ for the Riemann--Liouville integral based at
$z=1$, defined on $(z-1)^{\beta}C^\omega$ by analytic continuation in $\beta$; on powers it acts by
\begin{equation}\label{eq:RLpowers}
(z-1)^{\beta}\mapsto\frac{\Gamma(\beta+1)}{\Gamma(\beta+\nu+2)}(z-1)^{\beta+\nu+1}.
\end{equation}
For $\beta=\rho_1+k$ we have $\beta+\nu+2=k+1-e_1$, and $e_1\in i\R$, so the denominator never sits at a pole of
$\Gamma$.

\emph{Generic case.} If $\rho_1\notin\{-2,-1\}$, every factor in \eqref{eq:RLpowers} is finite and nonzero: for
$\Im\kappa_1\neq0$ the exponent $\rho_1$ is non-real, and for $\kappa_1=0$, $s=-2$ we have $\rho_1=2$. Put
$\tilde y=I^{\nu+1}y$. Then $\tilde y\neq0$ and $\tilde y\in(z-1)^{\rho_1+\nu+1}C^\omega=(z-1)^{-e_1}C^\omega$.

\emph{Threshold case} $\kappa_1=0$, $s=+2$, so $\rho_1=-2$. Write $y=(z-1)^{-2}(a_0+a_1(z-1))+y_{\mathrm{an}}$, with
$y_{\mathrm{an}}$ analytic at $1$. If $(a_0,a_1)\ne0$, let $\tilde y$ be the Pochhammer-contour version of the
transform. It annihilates analytic functions and maps $(z-1)^{-2},(z-1)^{-1}$ to nonzero multiples of
$(z-1)^{\nu-1},(z-1)^{\nu}$; here $\nu\notin\Z$ because $\kappa_3\ne0$ when $\kappa_1=0$. If $a_0=a_1=0$, then
$y$ is analytic at $1$, and we take $\tilde y=I^{\nu+1}y$, which is a convergent integral with nonzero factors in
\eqref{eq:RLpowers}. In both cases $\tilde y\ne0$ and $\tilde y\in(z-1)^{\nu-1}C^\omega=(z-1)^{-e_1}C^\omega$.

\emph{Analyticity at $z=t$.} Write $\tilde y(z)=(z-1)^{\nu+1}\int_0^1(1-\tau)^\nu y(1+(z-1)\tau)\,d\tau$, or the
corresponding Pochhammer integral. Since $y$ is analytic near $(1,t]$ and the singular points $0,\infty$ lie away
from the contour, $\tilde y$ is analytic near $z=t$. Finally, $\tilde y$ solves $\widetilde{\mathcal H}\tilde y=0$
because the boundary terms vanish: in the convergent case by the vanishing of the integrand at $w=1$, and for the
Pochhammer contour because it is closed.
\end{proof}

\begin{proposition}[The real dual]\label{prop:dual}
For real $\sigma\ne0$, after multiplication by an explicit nonvanishing gauge factor, the equation
$\widetilde{\mathcal H}\tilde y=0$ becomes
\begin{equation}\label{eq:dual}
-(\mu u')'-\frac{q(z)^2}{\mu}u+L(z)u=0,\qquad q(z)=q_1(z-z_0),\qquad L(z)=\frac{A}{z}+B-3z,
\end{equation}
with real coefficients, $q$ affine, and $q(1)$, $q(t)$ proportional (with a common sign) to
$\frac12(t-1)\Im e_1$ and $-\frac12t(t-1)\Im e_t$. If $w$ lies in the open window, then moreover:
\begin{itemize}
\item $q=q_1(z-z_0)$ with $q_1>0$ and $z_0\in(1,t)$, after replacing $u$ by $\bar u$ and $q$ by $-q$ if necessary;
\item $A=t\big(X_0^2-(X_1+X_t)^2\big)$ with $X_j=|\kappa_j+\kappa_3|$, and $A>0$;
\item $B=\lambda C_0+c_1$.
\end{itemize}
The transformed solution $u$ is outgoing for \eqref{eq:dual} in the sense of \cite[(3.4)]{Hin26}: near $z=1$ it is
$(z-1)^{-iq(1)/\mu'(1)}$ times a smooth function, and near $z=t$ it is $(t-z)^{-iq(t)/|\mu'(t)|}$ times a smooth
function.
\end{proposition}

\begin{proof}
For real $\sigma$ all exponents of $\widetilde{\mathcal H}$ at $0,1,t$ are purely imaginary, and the gauge factor $\prod_j(z-z_j)^{-\Re(\cdot)}$ removes the
real parts. The resulting operator is formally real, and writing it in the form \eqref{eq:dual} determines $q(1)$ and
$q(t)$ from the exponents at $z=1$ and $z=t$. We take $q$ to be linear. By Proposition~\ref{prop:window},
$q(1)q(t)<0$, so $q$ has a unique zero $z_0\in(1,t)$, and the sign of $q_1$ is fixed by replacing $u$ with $\bar u$
if necessary. The remaining claims are as follows.
\begin{itemize}
\item \emph{Form of $L$.} The function $L:=(q^2-\mu^2V)/\mu$ is then rational, and $L=n_0/z+n_1+n_2z$. Expanding at
$z=\infty$ with the exponents $\{1+2s,1+s-2\kappa_3\}$ gives $n_2=-3$, since the exponent difference there is $\pm4$.
\item \emph{The coefficient $n_0$.} Expanding at $z=0$ gives $n_0=t(X_0^2-(X_1+X_t)^2)$. In the window
$X_0-X_1-X_t=2\Im\kappa_3$, so $n_0=2t\,\Im\kappa_3\,(X_0+X_1+X_t)$. This is positive because $\Im\kappa_3>0$ for
$w>\Omega_3$, and $\Omega_3<\Omega_2<w$ since $|r_3|>r_2$.
\item \emph{The coefficient $n_1$.} The identity $n_1=\lambda C_0+c_1$ is an exact identity between rational
functions of $(r_0,r_1,r_2,\sigma,m,\lambda,a,s)$ after elimination by \eqref{eq:vieta} and \eqref{eq:residues}.
Its verification is computer algebra (Appendix~\ref{app:cert}, item~C1): the numerators of $n_1-(\lambda C_0+c_1)$
and $n_2+3$ are divisible by $s^2-4$.
\item \emph{Outgoing behavior.} This follows from Proposition~\ref{prop:transform}(ii) and the matching of exponents.
\qedhere
\end{itemize}
\end{proof}

\begin{lemma}[A priori bounds]\label{lem:floor}
Throughout the open window, $L'<0$ on $[1,t]$, $L(t)>0$ and $B\ge5.5t-1.25$.
\end{lemma}
\begin{proof}
First, $L'=-A/z^2-3<0$.

Next, by \eqref{eq:vieta}, $2(1-\alpha)C_0-(2t-1)$ is a positive multiple of $3(r_0+r_1)^2\ge0$, so
$C_0\ge\frac{2t-1}{2(1-\alpha)}$. This gives $c_1\ge4t-\frac12$.

Finally, the spin-weighted spheroidal eigenvalue satisfies $\lambda\ge\frac32$ on the window
(Proposition~\ref{prop:angular}). Hence $B=\lambda C_0+c_1\ge\frac32\cdot\frac{2t-1}{2}+4t-\frac12=5.5t-1.25$, and
$L(t)\ge B-3t>0$.
\end{proof}

\begin{proposition}[Angular bound]\label{prop:angular}
For $s=\pm2$ and $w=\sigma/m$ in the window, the eigenvalue satisfies $\lambda\ge\frac32$.
\end{proposition}

\begin{proof}
On the window $0<w$ (Proposition~\ref{prop:window}) and $aw\le1$, since $w_1$ is a positive-weight average of
$\Omega_1$ and $\Omega_3$, and each $\Omega_j<1/a$. Thus $\xi:=a\sigma$ satisfies $0\le m\xi\le m^2$. Also
$\alpha<7-4\sqrt3$ by Lemma~\ref{lem:alpha}. Write the angular operator as $H=H_0+P+\mathcal R$, where:
\begin{itemize}
\item $H_0$ is the operator at $\alpha=\xi=0$, with spectrum $\{\ell(\ell+1)-s(s+1)\}$;
\item $P$ is polynomial in $\cos\theta$ with coefficients linear in $\xi$;
\item $|\mathcal R|\le C\alpha$.
\end{itemize}

\emph{Case $|m|\ge3$.} The terms linear in $m\xi$ are nonnegative, and the remaining perturbation is bounded below by
$-6\alpha-4\alpha(1+\alpha)|m|$. This gives $\lambda\ge m^2+|m|-4.43-0.31|m|\ge6.6$.

\emph{Case $|m|\in\{1,2\}$.} We use a two-level bound. The lowest state of $H_0$ in the sector $(m,s)$ is the unique
$\ell=2$ state, and the next eigenvalue exceeds it by at least $6$. For a normalized trial state
$c\psi_0+\phi$ with $\phi\perp\psi_0$,
\[
\langle H(c\psi_0+\phi),c\psi_0+\phi\rangle\ge(2+\langle P\rangle-C\alpha)|c|^2+(8+\inf P-C\alpha)\|\phi\|^2
-2\langle P^2\rangle_c^{1/2}|c|\|\phi\|.
\]
The moments $\langle P\rangle$ and $\langle P^2\rangle$ are exact Beta integrals. The minimum over
$(\alpha,\xi)\in[0,0.0718]\times[0,|m|]$ is certified by interval arithmetic (Appendix~\ref{app:cert}, item~C2). The
certified minima are $1.957$, $1.665$, $1.920$ and $1.914$ for $(m,s)=(1,-2),(1,2),(2,-2),(2,2)$; all exceed $3/2$.
\end{proof}

\subsection{Currents}
On the dual equation \eqref{eq:dual} we use the currents of \cite[Lemma~4.1]{Hin26}, in the form with variable
potential from \cite[\S7]{Hin26}. Let $u_\pm=\exp(\mp i\int_{z_0}^zq/\mu)\,u$; these are smooth down to $z=t$ and
$z=1$, respectively. For a real function $f$, define
\begin{align*}
J^f[v]&=-\mu fq|v'|^2-\mu fq_1\Re(v'\bar v)+\big(Lfq+\tfrac12q_1\mu f'\big)|v|^2,\\
C^f&=\mu'fq-2\mu fq_1-\mu f'q,\qquad D^f=q(fL)'+\tfrac12q_1(\mu f')'.
\end{align*}
Then $\frac{d}{dz}J^f[u_\pm]=C^f|u_\pm'|^2+D^f|u_\pm|^2$. The proof is the integration by parts of
\cite[Lemma~4.1]{Hin26}, using $q'=q_1$.

\begin{lemma}[Horizon inequality]\label{lem:H}
For all $z_0\in[1,t]$ and $z\in(1,t)$, we have $H(z):=2\mu(z)+(z_0-z)\mu'(z)>0$.
\end{lemma}
\begin{proof}
$H$ is affine in $z_0$. At $z_0=1$ it equals $(z-1)\big((z-1)^2+t-1\big)>0$. At $z_0=t$ it equals
$(t-z)\big(t(t-1)-(t-z)^2\big)>0$.
\end{proof}

\begin{proposition}[Reduction to the event lemma]\label{prop:reduction}
Suppose there is a function $f\in C^2((1,z_0])\cap C^0([1,z_0])$ with the following properties:
\begin{enumerate}[label=(\alph*)]
\item $f(1)=0$, $f'(z)=o\big((z-1)^{-1}\big)$ as $z\to1$, and $f>0$ on $(1,z_0]$;
\item $f'(z_0)\ge0$;
\item with $s:=z_0-z$ and $h:=f'/f$, on $(1,z_0)$,
\begin{equation}\label{eq:CD}
\text{(C)}\quad s\,h\,\mu\le H,\qquad\qquad \text{(D)}\quad \tfrac12\big(\mu'h+\mu(h'+h^2)\big)\le s\,(Lh+L').
\end{equation}
\end{enumerate}
Then every outgoing solution of \eqref{eq:dual} vanishes.
\end{proposition}

\begin{proof}
\emph{Cosmological side.} On $[z_0,t]$ take $f_+=1$. Then $C^{f_+}=-q_1H\le0$ by Lemma~\ref{lem:H}, and
$D^{f_+}=qL'\le0$ because $q\ge0$ there. The limits of the current are computed as in \cite[Lemma~4.2]{Hin26}:
\[
\lim_{z\to t}J^{f_+}[u_+]=L(t)q(t)|B_t|^2,\qquad J^{f_+}[u_+](z_0)=-\mu(z_0)q_1\Re(u'\bar u)(z_0).
\]
Since $J^{f_+}$ is non-increasing and $L(t)q(t)>0$ by Lemma~\ref{lem:floor}, we get $\Re(u'\bar u)(z_0)\le0$, with
equality only if $B_t=0$.

\emph{Event side.} On $[1,z_0]$, write $q=-q_1s$. Then $C^f\le0$ is equivalent to (C), and $D^f\le0$ is equivalent to
$(\mu f')'\le2s(Lf)'$, which is (D). By (a), $\lim_{z\to1}J^f[u_-]=L(1)f(1)q(1)|B_1|^2=0$. Monotonicity gives
\[
J^f[u_-](z_0)=-\mu(z_0)q_1f(z_0)\Re(u'\bar u)(z_0)+\tfrac12q_1\mu(z_0)f'(z_0)|u(z_0)|^2\le0.
\]
So $\Re(u'\bar u)(z_0)\ge0$ by (b).

\emph{Conclusion.} Comparing the two sides forces $B_t=0$. At $z=t$ the indicial roots differ by a non-integer,
because $\Im e_t\ne0$ in the open window, so the regular-singular argument following \cite[(3.6)]{Hin26} gives
$u=0$.
\end{proof}

By Propositions~\ref{prop:transform}, \ref{prop:dual} and~\ref{prop:reduction}, the real-frequency part of
Theorem~\ref{thm:A} follows from the next result, which only involves the four numbers $(t,z_0,A,B)$.

\begin{theorem}[Event lemma]\label{thm:event}
For all $t>1$, $z_0\in(1,t)$, $A>0$ and $B\ge5.5t-1.25$, there exists $f$ satisfying (a)--(c) of
Proposition~\ref{prop:reduction}.
\end{theorem}

\section{Proof of the event lemma}\label{sec:event}

In terms of $f$, condition (C) reads $f'/f\le h_C:=H/(s\mu)=2/s+\mu'/\mu$. Condition (D) reads
$\widetilde Df:=(\mu f')'-2s(Lf)'\le0$. Put $x_0=z_0-1$ and
\[
p=\frac{2x_0L(1)}{t-1}.
\]
We distinguish three regimes:
\[
\text{R1: }p\le1;\qquad\text{R2: }p>1\text{ and }2x_0^2\Big(\frac{A}{z_0}+3\Big)<t-1;\qquad\text{R3: the rest.}
\]

\subsection{Regimes R1 and R2: an explicit solution}
\begin{proposition}\label{prop:R12}
In R1 and R2, the conclusion of Theorem~\ref{thm:event} holds.
\end{proposition}
\begin{proof}
Let $\tilde L=L$ in R1, and $\tilde L=L-\kappa$ with $\kappa=L(1)-\frac{t-1}{2x_0}>0$ in R2. Let $f$ solve
$(\mu f')'=2s(\tilde Lf)'$ with $f(1)=0$, that is,
\[
\mu f'=2s\tilde Lf+2\int_1^z\tilde Lf.
\]
At $z=1$ this is a regular-singular equation, and $f\sim(z-1)^{\tilde p}$ with
$\tilde p=2x_0\tilde L(1)/(t-1)\in(0,1]$. So (a) holds. In R2 we have $\tilde L(z_0)>0$, since
$L(1)-L(z_0)=Ax_0/z_0+3x_0$ and this is exactly the R2 condition. Thus $\tilde L>0$ on $[1,z_0]$, so $\mu f'>0$ and
$f$ is positive and increasing; in particular (b) holds.

For (D): $(\mu f')'=2s(\tilde Lf)'\le2s(Lf)'$, because $((L-\tilde L)f)'=\kappa f'\ge0$.

For (C): $\tilde L$ is decreasing, so
\[
s\mu f'\le2s\tilde L(1)f\,(s+z-1)=2sx_0\tilde L(1)f\le s(t-1)f.
\]
It remains to show that $H\ge s(t-1)$. With $z=1+x$,
\[
H-s(t-1)=x\big[2(1+x)(t-1-x)+s(2t-4-3x)\big].
\]
The bracket is affine in $s$. It is nonnegative at $s=0$, and at $s=t-1-x$ it equals $(t-1-x)(2t-2-x)\ge0$.
\end{proof}

\subsection{Regime R3: saturate (C), then solve (D)}
\begin{proposition}\label{prop:R3}
In R3, the conclusion of Theorem~\ref{thm:event} holds.
\end{proposition}

The construction follows \cite[\S4.3.3]{Hin26}. Since $h_C=(\log F_C)'$ with
\[
F_C:=\mu/s^2,
\]
we first take $f=F_C$, for which (C) holds with equality. Condition (D) then becomes
$d:=\widetilde DF_C/(2F_C)\le0$. A direct computation gives $\widetilde DF_C=N/s^4$, where in the variable
$S=z_0-z\in(0,x_0)$
\begin{equation}\label{eq:N}
N=6\mu_0^2-6\mu_0\mu_0'S+c_2S^2+c_3S^3+2(B+4t-15z_0+4)S^5+15S^6.
\end{equation}
Here $\mu_0=\mu(z_0)$, $\mu_0'=\mu'(z_0)$, and $c_2,c_3$ are affine in $(A,B)$ with explicit polynomial
coefficients in $(t,z_0)$. The coefficient of $S^4$ vanishes.

\begin{lemma}[One sign change]\label{lem:P1}
In R3, $N$ has exactly one zero $\hat S\in(0,x_0)$. Hence $d<0$ near $z=1$, $d>0$ near $z_0$, and $d$ changes sign
exactly once, at $\hat z=z_0-\hat S$.
\end{lemma}

\begin{proof}
Substitute $S=x_0v/(1+v)$, $v\in(0,\infty)$, and clear denominators. Let $T_0,\dots,T_6$ be the coefficients of the
resulting sextic in $v$.
\begin{itemize}
\item \emph{The first two coefficients.} $T_0=6\mu_0^2>0$, and $T_1=6z_0(z_0-t)x_0^2(3z_0^2+2z_0-4tz_0-t)>0$. Both
factors of $T_1$ are negative; the last factor is convex in $z_0$ and negative at $z_0=1$ and $z_0=t$.
\item \emph{The other coefficients.} For $j\ge2$, each $T_j$ is affine in $(A,B)$ with negative coefficients.
\item \emph{The parameter region.} In R3 the admissible $(A,B)$ form the quadrant
$\{A\ge a_0,\ B\ge5.5t-1.25\}$, where $a_0=\max\big(0,z_0((t-1)/(2x_0^2)-3)\big)$. This uses that the line
$A+B=3+(t-1)/(2x_0)$ coming from $p>1$ is redundant, which is a two-variable inequality.
\item \emph{Signs on the quadrant.} By the previous two items, it suffices to check the following at the corner of
the quadrant and on the triangle where $T_3\ge0$:
\[
T_4,T_5,T_6<0,\qquad\text{and}\qquad T_3\ge0\Rightarrow T_2>0.
\]
These reduce to finitely many polynomial inequalities in $(t,z_0)$. We certify them exactly in the variables
$(1/t,\zeta)$, with $\zeta=x_0/(t-1)$ (Appendix~\ref{app:cert}, item~C3).
\end{itemize}
Thus the sign pattern of $(T_0,\dots,T_6)$ has exactly one change. By Descartes' rule, $N$ has exactly one positive
root in $v$, and hence exactly one root in $S\in(0,x_0)$.
\end{proof}

\begin{proof}[Proof of Proposition~\ref{prop:R3}]
Define $f=F_C$ on $(1,\hat z]$, and on $[\hat z,z_0]$ let $f$ solve $\widetilde Df=0$ with $(f,f')(\hat z)=(F_C,F_C')(\hat z)$.
Since $d(\hat z)=0$, $f$ is $C^2$. Condition (a) holds because $F_C\sim(t-1)(z-1)/x_0^2$ as $z\to1$. Condition (D) holds on
$(1,\hat z]$ by Lemma~\ref{lem:P1}, and with equality on $[\hat z,z_0]$.

\emph{(C) on $[\hat z,z_0]$.} Let $p_*=\mu\exp\big(-\int2sL/\mu\big)>0$. Then
$\widetilde Df=(\mu/p_*)(p_*f')'-2sL'f$. The Wronskian $W=p_*(F_Cf'-F_C'f)$ satisfies
$W'=-(p_*/\mu)f\,\widetilde DF_C\le0$ wherever $f>0$, because $d\ge0$ on $[\hat z,z_0]$. Since $W(\hat z)=0$, we get
$f'/f\le F_C'/F_C=h_C$.

\emph{Positivity and the seam.} At a zero of $f'$ with $f>0$, $\mu f''=2sL'f<0$. Moreover
$f'(\hat z)=H(\hat z)/\hat s^3>0$. So it suffices to show $f'(z_0)>0$.

Suppose not, and let $z_1\in(\hat z,z_0]$ be the first zero of $f'$. The function $\Phi=\mu f'-2sLf$ satisfies
$\Phi'=2Lf\ge2L\hat f$ on $[\hat z,z_1]$, where $\hat f=f(\hat z)$. This yields
$\Psi(z_1)\le0$ for
\[
\Psi(z)=\hat H/\hat s-2\hat s\hat L+2\int_{\hat z}^zL+2s(z)L(z).
\]
But $\Psi'=2sL'<0$, so $\Psi(z_1)\ge\Psi(z_0)=\hat H/\hat s-2\int_{\hat z}^{z_0}s|L'|$. Using
$u-\log(1+u)\le u^2/2$,
\[
\int_z^{z_0}s|L'|=A\Big(\frac{z_0}{z}-1-\log\frac{z_0}{z}\Big)+\frac32s^2\le s^2\Big(\frac{A}{2z^2}+\frac32\Big),
\]
so it suffices that $K(\hat z)>0$, where $K=H-s^3(A/z^2+3)$.

Now $K$ is affine and decreasing in $A$, and $N$ is affine in $A$ and in $B$ with negative coefficients on
$(0,x_0)$. Hence the implication
\[
K\le0\ \Longrightarrow\ N(S;\max(A_K,0),5.5t-1.25)<0,
\]
where $A_K$ is the value of $A$ at which $K=0$, implies $K\le0\Rightarrow N<0$ for all admissible $(A,B)$. The
implication is a three-variable polynomial statement, which we certify exactly (Appendix~\ref{app:cert}, item~C4).
Since $N\ge0$ on $[0,\hat S]$ by Lemma~\ref{lem:P1}, it follows that $K(\hat z)>0$. Thus $f'>0$ on $[\hat z,z_0]$;
in particular $f>0$, and (b) holds.
\end{proof}

This completes the proof of Theorem~\ref{thm:event}, and with it Theorem~\ref{thm:A} for real $\sigma\ne0$.

\section{Zero frequency}\label{sec:zero}
\begin{proposition}\label{prop:zero}
There are no nontrivial outgoing modes at $\sigma=0$, for $s=\pm2$.
\end{proposition}
\begin{proof}
\emph{Case $m\ne0$.} Then $w=0<w_2$ lies outside the window, and Corollary~\ref{cor:outside} applies. Indeed,
$\Im\kappa_1$ and $\Im\kappa_3$ are nonzero because $am\ne0$.

\emph{Case $m=0$, $s=-2$.} The radial equation is $(\Delta_r^{-1}R')'+V\Delta_r^{-2}R=0$, with
$V=-2\Lambda r^2-2(1-\alpha)-\lambda$. Put $g=\Delta_r^2R''''$. The Teukolsky--Starobinsky identity at $K=0$ shows
that $g$ solves the same equation. Moreover $g$ vanishes like $\Delta_r^2$ at both horizons, so the energy identity
$\int\Delta_r^{-1}|g'|^2=\int V\Delta_r^{-2}|g|^2\le0$ gives $g=0$. Thus $R$ is a cubic polynomial. Inserting a cubic into the equation gives a $4\times4$ linear system for its
coefficients, with determinant $[\lambda^2+6(1-\alpha)\lambda+8-4\alpha+8\alpha^2]^2$. Both roots of the quadratic
are negative, and $\lambda\ge2-6\alpha>0$ at $m=\xi=0$ by the perturbation bound in the proof of
Proposition~\ref{prop:angular}. Hence the determinant is positive and $R=0$.

\emph{Case $m=0$, $s=+2$.} The function $\Delta_r^2R_{+2}$ solves the $s=-2$ equation, so we may repeat the argument.

The two identities used are exact symbolic identities (Appendix~\ref{app:cert}, item~C5).
\end{proof}

\section{The upper half plane}\label{sec:uhp}

\subsection{Fredholm theory}
\begin{lemma}[Regular frame]\label{lem:regframe}
Let $\chi\in C^\infty(\R)$ equal $1$ near $r_1$ and $0$ near $r_2$, and put $b(r)=\chi(r)+(1-\chi(r))/\Delta_r$ on
$(r_1,r_2)$. Let $(l_K,n_K,m_K)$ be the Kinnersley tetrad. Then
\[
l=b\,\Delta_r\,l_K,\qquad n=(b\,\Delta_r)^{-1}n_K,\qquad m=m_K
\]
extends to a stationary null tetrad that is smooth on $M=[r_1-\delta,r_2+\delta]\times S^2$, away from the usual phase
ambiguity of $m$ on the axis. In this tetrad, the GHP Teukolsky operator of spin $s=\pm2$ is a stationary second-order
operator with smooth coefficients on sections of a complex line bundle over $M$, with principal symbol
$\frac12G(\zeta,\zeta)$, where $G$ is the dual metric. Its outgoing modes are exactly the modes of
Lemma~\ref{lem:bc}.
\end{lemma}

\begin{proof}
Near $r_1$, in ingoing coordinates $(v,r,\tilde\varphi)$,
\[
\Delta_rl_K=2\Xi(r^2+a^2)\partial_v+2\Xi a\partial_{\tilde\varphi}+\Delta_r\partial_r,\qquad \Delta_r^{-1}n_K=-\partial_r/(2\varrho^2).
\]
Both are smooth and nonvanishing. Near $r_2$, in outgoing coordinates, $l_K=\partial_r+(\text{angular})$ and $n_K$ are
already smooth. The factor $b\Delta_r$ is smooth and positive on $(r_1,r_2)$, and equals $\Delta_r$ near $r_1$ and $1$
near $r_2$. The GHP operator is covariant, so its coefficients in this tetrad (spin coefficients and directional
derivatives) are smooth. Its principal part is $\Th\Th'-\eth\eth'$, whose symbol is
$\zeta(l)\zeta(n)-\zeta(m)\zeta(\bar m)=\frac12G(\zeta,\zeta)$. The last statement is the definition of outgoing
modes, together with Lemma~\ref{lem:bc}.
\end{proof}

Fix $m\in\Z$, and let $P_b(\sigma)$ be the spectral family of this operator on $\ker(\partial_{\varphi_*}-im)$, at
parameters $b=(\Lambda,M,a)$. The dynamical hypotheses of \cite[\S5.1]{HV18} concern only the principal symbol $G$
(radial sets at the horizons, and the global null-bicharacteristic flow). They were verified in the full
subextremal range by Petersen--Vasy \cite{PV25,PV24}; see \cite[Remark~2.12]{HPV25}. The subprincipal part enters
only through the threshold regularity at the radial sets. The trapping hypothesis of \cite[\S5.1]{HV18} is needed
only for high-energy estimates near the real axis, not for the Fredholm property \cite[Theorem~5.4 and the remark
in its proof]{HV18}. We therefore have:
\begin{enumerate}[label=(F\arabic*)]
\item $P_b(\sigma)$ is a holomorphic family of Fredholm operators of index $0$ on $\Im\sigma>-c$
\cite[Theorem~5.4]{HV18}. Its non-invertibility set is discrete, and it is the set of quasinormal mode (QNM)
frequencies.
\item There is $C_0$, uniform for $b$ near $b_0$, such that $P_b(\sigma)$ is invertible for $\Im\sigma>C_0$. This
follows from the global energy estimate on growing function spaces, which uses only hyperbolicity and holds
uniformly in $b$ \cite[\S5.1.2]{HV18}.
\item The non-invertibility set is closed in $(b,\sigma)$, and QNMs move continuously, with constant total rank in
small discs \cite[Proposition~5.11]{HV18}. That proof uses only the Fredholm estimates \cite[(5.26)]{HV18} on a
precompact set of $\sigma$, which hold uniformly in $b$.
\end{enumerate}
What is \emph{not} available for Teukolsky in the full subextremal range is the high-energy estimate bounding
$|\Re\sigma|$. We obtain it sector by sector from the gauge-fixed Einstein operator, for which it is proved in
\cite{HPV25}.

\subsection{Metric reconstruction}
We use the Geroch--Held--Penrose (GHP) formalism with $\Th'=\Delta-p\gamma-q\bar\gamma$ and
$\eth'=\bar\delta-p\alpha-q\bar\beta$ \cite{GHP73}. Let $\mathcal O$ be the GHP Teukolsky operator for $\Psi_0$, and
$\mathcal O'$ its prime. On any type D $\Lambda$-vacuum the Bianchi identities give $\Th\Psi_2=3\rho\Psi_2$,
$\eth\Psi_2=3\tau\Psi_2$ and their primes; there are no $\Lambda$-terms. With $P:=\Psi_2^{-4/3}$ (GHP type $(0,0)$),
and using the formal adjoint rules $\Th^\dagger=-(\Th-\rho-\bar\rho)$ and $\eth^\dagger=-(\eth-\tau-\bar\tau')$, a
direct computation gives
\begin{equation}\label{eq:adjoint}
\mathcal O'^\dagger\circ P=P\circ\mathcal O.
\end{equation}
For instance, the radial part is
\[
(\Th-\bar\rho)(\Th'+3\rho')(PX)=P(\Th-4\rho-\bar\rho)(\Th'-\rho')X.
\]
Wald's argument \cite{Wal78}, which uses only the decoupling identity $\mathcal S'\mathcal E=\mathcal O'\mathcal T'$
and the self-adjointness of the linearized Einstein operator $\mathcal E$, then gives the following.

\begin{lemma}[Reconstruction]\label{lem:recon}
If $\mathcal O\Psi=0$, then $h_\Psi:=\mathcal S'^\dagger(P\Psi)$ solves $\mathcal E h_\Psi=0$. Moreover, the complex
conjugate $\tilde h:=\overline{h_\Psi}$ satisfies $\Psi_4[\tilde h]=\frac12(\Th')^4(\bar P\bar\Psi)$.
\end{lemma}

\begin{proof}
The first claim is Wald's argument combined with \eqref{eq:adjoint}. For the second: $h_\Psi$ has components only
along $l\otimes l$, $l\otimes m$ and $m\otimes m$. By GHP type counting, $\delta\Psi_4$ involves only the components
$h_{nn}$, $h_{n\bar m}$ and $h_{\bar m\bar m}$. Hence $\Psi_4[h_\Psi]=0$, and $\Psi_4[\tilde h]$ comes from
$\tilde h_{\bar m\bar m}$ alone. It involves only $\Th'$-derivatives and the
coefficients $\rho',\bar\rho'$, and its evaluation uses only the Ricci identities $\Th'\rho'=\rho'^2$ and
$\Th'\bar\rho'=\bar\rho'^2$. In particular, no commutator enters, so the $\Lambda$-terms of the GHP commutators play
no role. The algebra is therefore literally that of the Kerr case \cite{Wal78,KSFK03}.
\end{proof}

\subsection{The kernel lemma}
Let $\kappa_2^\sharp=|\Delta_r'(r_2)|/(2\Xi\rho_2)$ be the normalized surface gravity of the cosmological horizon,
and define the \emph{cosmological ladder}
\[
\mathcal Z_m=\big\{m\Omega_2+iN\kappa_2^\sharp:\ N\in\{1,2,3,\dots\}\big\}.
\]

\begin{lemma}[Kernel lemma]\label{lem:kernel}
Let $\Im\sigma>0$ and $\sigma\notin\mathcal Z_m$. If $F=e^{-i\sigma t_*+im\varphi_*}v$ is a mode of GHP type
$(0,4)$ or $(4,0)$, regular at the future cosmological horizon, with $(\Th')^4F=0$, then $F=0$.
\end{lemma}

\begin{proof}
In the tetrad of Lemma~\ref{lem:regframe}, $\Th'$ differentiates only along $n$, which has no
$\partial_\theta$-component. So for each fixed $\theta$, $(\Th')^4F=0$ is a fourth-order linear ODE in $r$, the
composition of four first-order operators
\[
n(\cdot)-(p_i+q_i)\gamma,\qquad p_i+q_i\in\{4,2,0,-2\}.
\]
On modes, $n=-(\Delta_r/2\varrho^2)\big(\partial_r+iK/\Delta_r\big)$ in the Kinnersley normalization, and
$\gamma=\mu+\Delta_r'/(4\varrho^2)$ with $\mu$ vanishing like $\Delta_r$. Hence each first-order
factor has kernel spanned by $\exp(-i\int K/\Delta_r)$ times $\Delta_r^{k}$ times an analytic function, with
$k\in\frac12(p_i+q_i)+\Z\subset\Z$. Every local solution near $r_2$ is therefore a finite sum of terms
$(r_2-r)^{-\kappa_2+k}\log^j(r_2-r)$ times analytic functions, with $k\in\Z$.

Regularity at the future cosmological horizon is the branch $(r_2-r)^{\kappa_2}C^\infty$ (Lemma~\ref{lem:bc}). So a
nonzero $F$ requires $2\kappa_2\in\Z$. For $\Im\sigma>0$ we have
$\Re(2\kappa_2)=2\Xi\rho_2\Im\sigma/|\Delta_r'(r_2)|>0$ and $\Im(2\kappa_2)=2\Xi(\rho_2\Re\sigma-am)/\Delta_r'(r_2)$.
Hence $\Re\sigma=m\Omega_2$ and $\Im\sigma=N\kappa_2^\sharp$ with $N=2\kappa_2\in\{1,2,3,\dots\}$. Finally, $F$ vanishes identically once it vanishes near $r_2$, by uniqueness for the ODE.
\end{proof}

\begin{proposition}[Transfer]\label{prop:transfer}
Let $\Psi$ be an outgoing $s=+2$ Teukolsky mode at frequency $\sigma$ in the sector $m$, with $\Im\sigma>0$ and
$\sigma\notin\mathcal Z_m$. Then $\tilde h$ from Lemma~\ref{lem:recon} is a smooth outgoing mode of the linearized
Einstein equation, at frequency $-\bar\sigma$ in the sector $-m$, with $\Psi_4[\tilde h]\ne0$. The case $s=-2$ is
the same with the roles of $l$ and $n$, and of the two horizons, exchanged: one uses the ingoing gauge and the
ladder $\mathcal Z_m^e$ of the event horizon.
\end{proposition}

\begin{proof}
\emph{Smoothness.} $\mathcal S'^\dagger$ is GHP-covariant. Near each future horizon it may therefore be evaluated in
a principal tetrad that is smooth there. In such a tetrad, outgoing modes are smooth, so $\tilde h$ is smooth.

\emph{Nonvanishing of $\Psi_4$.} The function $\bar P\bar\Psi$ is regular at both horizons, and its frequency
$-\bar\sigma$ lies off $\mathcal Z_{-m}=-\overline{\mathcal Z_m}$. By Lemma~\ref{lem:kernel},
$(\Th')^4(\bar P\bar\Psi)\ne0$.
\end{proof}

\subsection{Compactness, sector by sector}
\begin{lemma}\label{lem:compact}
Let $b_0$ be subextremal and $m\in\Z$. There are a neighborhood $U\ni b_0$ and $R_m<\infty$ such that, for all
$b\in U$, every Teukolsky QNM of spin $\pm2$ in the sector $m$ with $\Im\sigma\ge0$ satisfies $|\sigma|\le R_m$.
\end{lemma}

\begin{proof}
Let $L_b$ be the linearized gauge-fixed Einstein operator of \cite[(23)]{HPV25}, with the constraint damping fixed at
$b_0$. By \cite[Remark~2.12 and \S2.5]{HPV25}, the operators $L_b$, $b\in U$, form a continuous family of stationary,
principally scalar operators satisfying the assumptions of \cite[\S5.1]{HV18}, including the trapping condition
(by \cite[Theorem~3.1]{HPV25}). By \cite[\S5.1.2]{HV18}, these assumptions then hold uniformly near $b_0$, and the
high-energy estimates \cite[(5.7), (5.8)]{HV18} hold with constants uniform in $b\in U$. In particular there is $R_0$
such that every QNM of $L_b$ with $\Im\ge0$ has $|\sigma_T|\le R_0$ for all $b\in U$. Here $\sigma_T$
is the $T$-frequency, which differs from the $\partial_t$-frequency by $m'\Omega_T$ on the sector $m'$.

Let $\Psi$ be a QNM in the sector $m$ with $\Im\sigma\ge0$. If $\Im\sigma=0$ there is nothing to prove, by
Theorem~\ref{thm:A} on the real axis. By (F2), $\Im\sigma\le C_0$, so the ladder points in the strip
$0<\Im\sigma\le C_0$ satisfy $|\sigma|\le|m|\Omega_2+C_0$, or $|\sigma|\le|m|\Omega_1+C_0$ for the event-horizon ladder. If $\sigma$ is not a ladder point,
Proposition~\ref{prop:transfer} gives $\tilde h$ at frequency $-\bar\sigma$ with $\Psi_4[\tilde h]\ne0$.

Next we gauge-fix $\tilde h$: put $h'=\tilde h+\delta^*\omega$, where $\omega$ solves the linearized gauge condition
$D\Upsilon(\delta^*\omega)=-D\Upsilon(\tilde h)$ at frequency $-\bar\sigma$. The operator $\omega\mapsto D\Upsilon(\delta^*\omega)$
is a principally scalar wave operator on 1-forms, so it is Fredholm of index $0$ by (F1) in its version for
1-forms.
\begin{itemize}
\item If it is invertible at $-\bar\sigma$, then $h'$ is a QNM of $L_b$, and $h'\ne0$ since $\Psi_4[h']=\Psi_4[\tilde h]\ne0$.
\item If it is not invertible, there is $\omega\ne0$ with $D\Upsilon(\delta^*\omega)=0$. Then
$L_b(\delta^*\omega)=0$, and $\delta^*\omega\ne0$ because KdS has no Killing modes with $\Im\ne0$. So
$\delta^*\omega$ is a QNM of $L_b$ at frequency $-\bar\sigma$.
\end{itemize}
In either case $L_b$ has a QNM at $T$-frequency $-\bar\sigma-(-m)\Omega_T$, so
$|\sigma|\le R_0+|m|\Omega_T$. For $s=-2$ the event-horizon ladder enters instead. Hence we can take
$R_m=\max(R_0+|m|\Omega_T,\ |m|\Omega_1+C_0,\ |m|\Omega_2+C_0)$.
\end{proof}

\subsection{The base case}
\begin{lemma}[Schwarzschild--de Sitter]\label{lem:sds}
For all but countably many values of $\Lambda M^2\in(0,1/9)$, SdS has no Teukolsky QNMs of spin $\pm2$ with
$\Im\sigma>0$.
\end{lemma}

\begin{proof}
Let $a=0$. The angular operator then does not depend on $\sigma$ or $\Lambda$, so modes separate into spin-weighted
spherical harmonics for every $\sigma\in\C$, with a discrete real spectrum $\{\lambda_\ell\}$ independent of $\Lambda$. Two
exact computations hold on separated modes, in the Kinnersley tetrad with $D^\dagger=\partial_r+i\sigma r^2/\Delta_r$
and $D=\partial_r-i\sigma r^2/\Delta_r$ (Appendix~\ref{app:cert}, item~C8):
\begin{equation}\label{eq:sdsTS}
(\Th')^4(r^4R)=\frac{1}{16r^4}\,\Delta_r^2(D^\dagger)^4\Delta_r^2R,\qquad
D^4\Delta_r^2(D^\dagger)^4\Delta_r^2R=\mathcal C_\ell(\sigma)R,
\end{equation}
where $\mathcal C_\ell(\sigma)=[(\lambda_\ell+4)(\lambda_\ell+6)]^2+144M^2\sigma^2$, the second identity holds for every solution
$R$ of the $s=+2$ radial equation, and $\Lambda$ drops out. Since $P\propto r^4$ at $a=0$, the first identity shows
that $(\Th')^4(P\Psi)$ vanishes only if $\Delta_r^2(D^\dagger)^4\Delta_r^2R=0$, hence only if $\mathcal C_\ell(\sigma)=0$, that is,
$\sigma=\sigma_\ell:=i(\lambda_\ell+4)(\lambda_\ell+6)/(12M)$.

By Lemma~\ref{lem:kernel} (with $\Omega_2=0$), this can happen for a mode regular at both horizons only if
$\sigma_\ell=iN\kappa_2^\sharp$ for some integer $N\ge1$, that is, only if $M\kappa_2^\sharp\in\{(\lambda_\ell+4)(\lambda_\ell+6)/(12N)\}$.
The right-hand side is a countable set independent of $\Lambda$. On the other hand, $M\kappa_2^\sharp$ is a nonconstant real-analytic function of
$\Lambda M^2\in(0,1/9)$, so this happens only for countably many values of $\Lambda M^2$.

Off this set, Lemma~\ref{lem:recon} and the argument of Proposition~\ref{prop:transfer} produce, from a QNM with
$\Im\sigma>0$, a linearized Einstein mode with $\Im>0$ and $\Psi_4\ne0$. By \cite[\S7]{HV18}, every such mode on SdS
is pure gauge, and pure gauge modes have $\Psi_4=0$. This is a contradiction. The case $s=-2$ is the same, using the
ingoing gauge, $\Th^4$, and the event horizon.
\end{proof}

\subsection{Conclusion}
\begin{proof}[Proof of Theorem~\ref{thm:A} for $\Im\sigma>0$]
Fix a subextremal $b_1$ and $m\in\Z$. Choose an SdS point $b_*$ as in Lemma~\ref{lem:sds}, and a continuous path
$p:[0,1]\to\{\text{subextremal parameters}\}$ with $p(0)=b_*$ and $p(1)=b_1$. For example, scale $r_0$ to $0$ with
$r_1,r_2$ fixed, as in \cite[\S6]{Hin26}, and then move within the SdS family. Let $S$ be the set of $s\in[0,1]$ for
which the sector $m$ has a QNM with $\Im\sigma>0$.
\begin{itemize}
\item $S$ is open, by (F3).
\item $S$ is closed. Let $s_j\to s_0$ with $s_j\in S$ and QNMs $\sigma_j$. By Lemma~\ref{lem:compact} and (F2), the
$\sigma_j$ lie in the compact set $\{0\le\Im\sigma\le C_0,\ |\sigma|\le R_m\}$. A subsequence converges to some
$\sigma_0$, which is a QNM at $p(s_0)$ by (F3). Now $\sigma_0\ne0$, because $P(0)$ is invertible by
Proposition~\ref{prop:zero} and invertibility is open. And $\sigma_0\notin\R$ by \S\ref{sec:event} and
Corollary~\ref{cor:outside}. Hence $\Im\sigma_0>0$.
\item $0\notin S$, by Lemma~\ref{lem:sds}.
\end{itemize}
Hence $S=\emptyset$. Generalized modes reduce to modes, since their top coefficient is a mode.
\end{proof}

\section{Ungauged linearized Einstein mode stability}\label{sec:B}

Throughout this section $a\ne0$, and $\Omega_T=a/(r_2^2+a^2)$, so that $T=\partial_t+\Omega_T\partial_\varphi$ in
Boyer--Lindquist coordinates. Let $h$ be as in Theorem~\ref{thm:B}. Its Fourier components in $\varphi_*$ are
generalized $(\partial_t,\partial_\varphi)$-modes of type $(\sigma+m\Omega_T,m)$, whose imaginary part is $\Im\sigma\ge0$. The linearized
Teukolsky scalars $\delta\Psi_0$ and $\delta\Psi_4$ are gauge and tetrad invariant on type D, and they are outgoing
generalized Teukolsky modes. By Theorem~\ref{thm:A}, applied to each Fourier component of $h$ and of $\bar h$, the scalars $\delta\Psi_0$ and
$\delta\Psi_4$ vanish on the exterior $r_1<r<r_2$. The next lemma extends this to all of $M$:
\[
\delta\Psi_0=\delta\Psi_4=0\quad\text{on }M.
\]

\begin{lemma}[Collars]\label{lem:collar}
Let $w$ be a smooth generalized mode on $M$ of a stationary, principally scalar operator with smooth coefficients
and principal symbol $\frac12G$, for instance the Teukolsky operator of Lemma~\ref{lem:regframe}. If $w=0$ on
$r_1<r<r_2$, then $w=0$ on $M$.
\end{lemma}

\begin{proof}
Consider the collar $r_1-\delta<r<r_1$. There $dr$ is timelike, and $r$ decreases towards the future. Let $p$ lie in
the collar. Every past-directed causal curve from $p$ has strictly increasing $r$. The set
$D_p=J^-(p)\cap\{r\le r_1\}$ is compact, since in ingoing coordinates the past cone of $p$ meets the event horizon
in a compact set. Its boundary consists of a compact piece of the event horizon and a piece of the past light cone
of $p$. Since $w$ is smooth and vanishes for $r>r_1$, it vanishes to infinite order on the horizon. The standard
energy estimate on $D_p$ involves fluxes through a null boundary that are tangential and hence vanish, and fluxes
through the cone of $p$ that are nonnegative \cite[\S23.2]{Hor85}. It gives $w=0$ near $p$. The collar
$r_2<r<r_2+\delta$ is treated in the same way, with time reversed relative to $r$ and the cosmological horizon in
place of the event horizon.
\end{proof}

\subsection{Linearized Goldberg--Sachs}
\begin{lemma}\label{lem:GS}
After an infinitesimal null rotation of the tetrad, $\delta\Psi_1=\delta\Psi_3=0$, and then
$\delta\kappa=\delta\sigma=\delta\kappa'=\delta\sigma'=0$.
\end{lemma}
\begin{proof}
To first order, a null rotation about $n$ with parameter $c$ shifts $\delta\Psi_1$ by $3c\Psi_2$ and fixes
$\delta\Psi_0$ and $\delta\Psi_2$; the analogous statement holds about $l$. Since $\Psi_2\ne0$, the choice is
algebraic. The $\Lambda$-vacuum Bianchi identities contain no $\Lambda$-terms. Linearizing
\[
\Th\Psi_1-\eth'\Psi_0=-\tau'\Psi_0+4\rho\Psi_1-3\kappa\Psi_2,\qquad
\Th'\Psi_0-\eth\Psi_1=\rho'\Psi_0-4\tau\Psi_1+3\sigma\Psi_2
\]
gives $\delta\kappa=\delta\sigma=0$; priming gives the rest.
\end{proof}

All perturbed quantities below are local linear functions of $h$. They are boost-spin invariant and commute with
$\cL_{\partial_t}$ and $\cL_{\partial_\varphi}$.

\subsection{The Killing-spinor field}
On a type D $\Lambda$-vacuum with $\kappa=\sigma=\kappa'=\sigma'=0$, the following identities hold, together with
their primes:
\begin{gather*}
\Th\rho=\rho^2,\quad\eth\rho=(\rho-\bar\rho)\tau,\quad\Th\tau=(\tau-\bar\tau')\rho,\quad\eth\tau=\tau^2,\\
\Th\rho'-\Th'\rho=\tau\bar\tau-\tau'\bar\tau',\quad\eth\tau'-\eth'\tau=\rho\bar\rho'-\bar\rho\rho',\quad
\Th\tau'=\eth'\rho,\quad\Th'\tau=\eth\rho'.
\end{gather*}
They follow from the GHP Ricci identities and commutators. We verified them exactly on KdS, and on the accelerating NUT
Pleba\'nski--Demia\'nski metric, which has $\tau\ne-\bar\tau'$ (Appendix~\ref{app:cert}, item~C6).

Let $\phi=\Psi_2^{-1/3}$ (normalized so that $\phi=r-ia\cos\theta$ on KdS), and put
\[
A=\phi\rho,\quad B=-\phi\rho',\quad C=\phi\tau,\quad D=-\phi\tau',\qquad \xi=Bl+An-Dm-C\bar m.
\]

\begin{lemma}\label{lem:xi}
The ten Killing equations for $\xi$ are algebraic consequences of the identities above. On KdS,
$\xi=-\Xi\,\partial_t$.
\end{lemma}

Linearizing Lemma~\ref{lem:xi} gives $\cL_\xi h+\cL_{\delta\xi}g=0$, where $\delta\xi=\delta\xi[h]$.

\begin{lemma}[Reality]\label{lem:reality}
Put $J[h]=\big(\delta\xi[h]-\overline{\delta\xi[\bar h]}\big)/(2i)$. Then $\cL_{J[h]}g=0$. If $(\sigma,m)\ne(0,0)$,
then $J[h]=0$.
\end{lemma}
\begin{proof}
For real $h$, $J[h]=\Im\delta\xi$, and the imaginary part of the linearized Killing equation reads $\cL_Jg=0$;
by complex linearity this holds for all $h$.

The local Killing algebra of KdS with $a\neq0$ is $\mathrm{span}\{\partial_t,\partial_\varphi\}$ on every open set. Indeed,
a Killing field preserves $\Psi_2$, hence $r$ and $a\cos\theta$. It is therefore tangent to the $(t,\varphi)$-orbits,
and the isotropy is trivial because it must fix $\nabla r$, $\nabla(a\cos\theta)$ and the principal null directions.

Moreover, $J$ is complex linear and commutes with $\cL_{\partial_t}$ and $\cL_{\partial_\varphi}$. For a mode of type
$(\sigma,m)\ne(0,0)$, it is therefore a constant combination of $\partial_t,\partial_\varphi$ of that type, hence zero.
\end{proof}

\begin{remark}\label{rem:Jvanish}
The same argument shows $J[h]=0$ in two further cases. If $(\cL_T+i\sigma)h=0$ with $\sigma\ne0$, then
$\cL_TJ[h]=-i\sigma J[h]$, while $\cL_T$ annihilates every Killing field of KdS. If $\cL_Th=0$ and $h$ has vanishing
$\varphi$-average, then $J[h]$ is a constant combination of $\partial_t,\partial_\varphi$, hence equal to its own
$\varphi$-average, which is $J$ of the $\varphi$-average of $h$, namely $0$.
\end{remark}

Consequently $\phi\rho=\bar\phi\bar\rho$, $\phi\rho'=\bar\phi\bar\rho'$ and $\phi\tau=-\bar\phi\bar\tau'$ hold to
first order.

\subsection{Carter's Killing field}
Let $R=\Re\phi$ and $I=\Im\phi$. These are the Carter coordinates $r$ and $-a\cos\theta$. Reality gives
\[
\Th R=-A,\quad \Th'R=B,\quad \eth R=0,\quad \Th I=\Th'I=0,\quad \eth I=iC.
\]
Define
\[
\zeta=I^2(Bl+An)+R^2(Dm+C\bar m).
\]

\begin{lemma}\label{lem:zeta}
To first order, $\zeta$ is Killing. On KdS, $\zeta=-\Xi(a^2\partial_t+a\partial_\varphi)$.
\end{lemma}
\begin{proof}
The $ll$, $nn$, $mm$ and $\bar m\bar m$ components are immediate from those of $\xi$.
\begin{itemize}
\item The $ln$ component equals $-|\phi|^2\phi(\tau\bar\tau-\tau'\bar\tau')=0$, since by reality
$\tau\bar\tau=\tau'\bar\tau'$.
\item The $m\bar m$ component equals $|\phi|^2\phi(\rho\bar\rho'-\rho'\bar\rho)=0$, again by reality.
\item The $lm$ component equals $2AC\phi+|\phi|^2(\eth A+\bar\tau'A)=0$, using $\eth A=-\phi\bar\rho\tau$.
\item The $nm$ component equals $-2BC\bar\phi+|\phi|^2(\eth B+\tau B)$. Differentiating the first-order reality
relation $\phi\rho'=\bar\phi\bar\rho'$ along $m$ gives $\eth B+\tau B=2\tau B$, and the sum vanishes.
\end{itemize}
The remaining components are complex conjugates. The value on KdS is an exact computation (item~C6).
\end{proof}

\subsection{Pure gauge modes}
On KdS, $\partial_t=-\xi/\Xi$ and $\partial_\varphi=(a^2\xi-\zeta)/(\Xi a)$. Hence there are constants $\alpha_T,\beta_T$
and $\alpha_\varphi,\beta_\varphi$ with
\[
T=\alpha_T\xi+\beta_T\zeta,\qquad \partial_\varphi=\alpha_\varphi\xi+\beta_\varphi\zeta.
\]
Whenever $J[h]=0$, Lemmas~\ref{lem:xi} and~\ref{lem:zeta} give $\cL_\xi h=-\cL_{\delta\xi}g$ and
$\cL_\zeta h=-\cL_{\delta\zeta}g$, and therefore
\begin{equation}\label{eq:LTh}
\cL_Th=-\cL_{Y_T}g,\qquad \cL_{\partial_\varphi}h=-\cL_{Y_\varphi}g,\qquad
Y_\bullet=\alpha_\bullet\delta\xi+\beta_\bullet\delta\zeta.
\end{equation}
The vector fields $Y_T$ and $Y_\varphi$ are local linear expressions in $h$. They are smooth across the horizons:
$\xi$, $\zeta$, $\delta\xi$ and $\delta\zeta$ are boost-spin invariant constructions from the perturbed Weyl spinor,
so they can be computed in a horizon-regular principal tetrad, and $\Psi_2\ne0$ on $r>0$.

\begin{proposition}\label{prop:puregauge}
\begin{enumerate}[label=(\roman*)]
\item If $(\cL_T+i\sigma)h=0$ with $\sigma\ne0$, then $h=\cL_Xg$ with $X=Y_T/(i\sigma)$.
\item If $\cL_Th=0$ and $h$ has vanishing $\varphi$-average, then $h=\cL_Xg$, where $X$ is the unique
$\varphi$-antiderivative of $-Y_\varphi$ with vanishing $\varphi$-average.
\end{enumerate}
In both cases $X$ is smooth, and $(\cL_T+i\sigma)X=0$.
\end{proposition}

\begin{proof}
In both cases $J[h]=0$ by Remark~\ref{rem:Jvanish}, so \eqref{eq:LTh} holds.

(i) Here $-i\sigma h=\cL_Th=-\cL_{Y_T}g$.

(ii) The field $Y_\varphi$ has vanishing $\varphi$-average, because $Y_\varphi$ commutes with rotations and $h$ has
vanishing average. So its mean-zero $\varphi$-antiderivative $X$ (with $[\partial_\varphi,X]=-Y_\varphi$) is smooth.
Then $\cL_{\partial_\varphi}(h-\cL_Xg)=-\cL_{Y_\varphi}g+\cL_{Y_\varphi}g=0$. So $h-\cL_Xg$ is axisymmetric with
vanishing average, hence zero.
\end{proof}

\subsection{The stationary axisymmetric sector}
\begin{proposition}\label{prop:stationary}
If $\cL_{\partial_t}h=\cL_{\partial_\varphi}h=0$ and $\delta\Psi_0=\delta\Psi_4=0$, then
$h=c_M\partial_Mg+c_a\partial_ag+\cL_Xg$, with $X$ affine in $(t,\varphi)$.
\end{proposition}

\begin{proof}
\emph{Step 1: restoring reality.} Here $J[h]=A'\partial_t+B'\partial_\varphi$ need not vanish. A constant phase change of
$\phi$ removes $A'$. Subtracting a multiple of the tangent to the Pleba\'nski--Demia\'nski family in the direction of
the acceleration removes $B'$: this is the only Pleba\'nski--Demia\'nski tangent that makes $\xi$ complex, and its
$J$ has nonzero $\partial_\varphi$-part (exact nonlinear check, item~C6). After this, reality holds, and $\xi$ and
$\zeta$ are first-order Killing fields. They commute, since $\delta\xi$ and $\delta\zeta$ are background Killing
fields.

\emph{Step 2: Carter's ansatz.} One has $\nabla R\perp\nabla I$. Put $\Delta_R:=-2AB|\phi|^2$ and
$\Delta_I:=-2CD|\phi|^2$. Then $\Delta_R$ depends on $R$ alone and $\Delta_I$ on $I$ alone, since
$\eth(AB|\phi|^2)=0$ by reality and $\xi(\Delta_R)=0$. All the inner products of $\xi$, $\zeta$, $\nabla R$ and $\nabla I$ are
functions of $(R,I,\Delta_R,\Delta_I)$. In coordinates adapted to $(\xi,\zeta,R,I)$, the perturbed metric is therefore,
to first order, Carter's separable ansatz
\[
\frac{S}{\Delta_R}dR^2+\frac{S}{\Delta_I}dI^2+\frac{\Delta_I(d\tilde t+R^2d\psi)^2-\Delta_R(d\tilde t-I^2d\psi)^2}{S},\qquad S=R^2+I^2.
\]

\emph{Step 3: the field equations.} On this ansatz the $\Lambda$-vacuum equations are linear in $(\Delta_R,\Delta_I)$,
with general solution
\[
\Delta_R=-\tfrac\Lambda3R^4+\varepsilon R^2-2\mathfrak mR+\gamma,\qquad \Delta_I=-\tfrac\Lambda3I^4-\varepsilon I^2+2\mathfrak nI+\gamma
\]
(exact computation, item~C7). Modulo the rescaling gauge, the variations are $\delta M$, $\delta a$ and the NUT
variation $\delta\mathfrak n$, together with the acceleration removed in Step 1.

\emph{Step 4: axis regularity.} Elementary flatness at the two axes gives two linear conditions on
$(\delta\mathfrak n,\delta\alpha_{\mathrm{acc}})$ with determinant $-64Mp_0^4/c\ne0$, so both vanish. The gauge
vector $X$ is the generator of the change from the background to the perturbed adapted coordinates, together with
the rescalings of $\xi$ and $\zeta$.

\emph{Step 5: the gauge vector on all of $M$.} Here $\partial_Mg$ and $\partial_ag$ denote the parameter derivatives
of the KdS family in the regular coordinates $(t_*,r,\theta,\varphi_*)$, as in \cite{HPV25}.

The parameters $c_M,c_a$ are determined by local invariants, for example $\delta(\Psi_2\phi^3)$, which is smooth on $M$
and locally constant on each of the three regions $M\setminus\{r=r_1,r_2\}$. Hence they are the same on all three
regions. By Lemma~\ref{lem:collar}, Steps 1--4 apply verbatim on each of the three regions, where
$\Psi_2\ne0$ and $\Delta_R\ne0$. So $k:=h-c_M\partial_Mg-c_a\partial_ag$ is smooth on $M$, and $k=\cL_{X^{(i)}}g$ on
the $i$-th region with $X^{(i)}$ smooth there.

We now glue these. Write $X=X^r\partial_r+X^\theta\partial_\theta+Y$ with $Y=Y^t\partial_{t_*}+Y^\varphi\partial_{\varphi_*}$.
\begin{itemize}
\item \emph{The $(r,\theta)$-components.} Since $R=\Re\phi$ and $I=\Im\phi$ are natural functions,
$X(R)=\delta R[k]$ and $X(I)=\delta I[k]$. These are local linear expressions in $k$, smooth on $M$. Hence
$X^r$ and $X^\theta$ are given by one formula on all three regions, and they are smooth on $M$.
\item \emph{The remaining components.} Let $k_1=k-\cL_{X^r\partial_r+X^\theta\partial_\theta}g$, which is smooth on $M$. On
each region, $Y^\bullet=c^\bullet t_*+f^\bullet(r,\theta)$ with constants $c^\bullet$. Since $\partial_{t_*},\partial_{\varphi_*}$ are
Killing,
\[
k_1=\cL_Yg=2\,dY^t\odot K_t^\flat+2\,dY^\varphi\odot K_\varphi^\flat,\qquad K_t=\partial_{t_*},\ K_\varphi=\partial_{\varphi_*}.
\]
\emph{Claim:} at every point off the axis, the linear map $(dY^t,dY^\varphi)\mapsto\cL_Yg$, restricted to pairs of
covectors annihilating $\partial_{\varphi_*}$, is injective. Indeed, $K_t^\flat,K_\varphi^\flat$ are linearly independent, and
$\alpha\odot K_t^\flat+\beta\odot K_\varphi^\flat=0$ forces $\alpha=yK_\varphi^\flat$ and $\beta=-yK_t^\flat$. Then
$\alpha(\partial_{\varphi_*})=y\,g_{\varphi\varphi}=0$ gives $y=0$.

Hence, off the axis, $(dY^t,dY^\varphi)$, and with it $(c^\bullet,df^\bullet)$, is given by one smooth formula in $k_1$, the same on
all three regions. Integrating, $Y$ is smooth on $M$ minus the axis, up to additive constants, which are Killing. Near the
axis, $Y$ is smooth on each region. It extends smoothly across the codimension-three set where the axis meets the
horizons by Killing transport: $(Y,\nabla Y)$ satisfies a linear ODE with coefficients smooth in $(g,k_1)$ along
any curve.
\end{itemize}
So $X$ is smooth on $M$. Then $k-\cL_Xg$ is smooth on $M$ and vanishes on each of the three regions, hence on $M$.
The same argument applies to the $t_*$-linear parts that arise in the generalized-mode induction below.
\end{proof}

\subsection{Generalized modes and proof of Theorem~\ref{thm:B}}
\begin{proof}[Proof of Theorem~\ref{thm:B}]
\emph{Part (1).} Let $(\cL_T+i\sigma)^kh=0$ with $\sigma\ne0$. We argue by induction on $k$; the case $k=1$ is
Proposition~\ref{prop:puregauge}(i). For $k\ge2$, the top coefficient $h_1=(\cL_T+i\sigma)^{k-1}h$ is a mode, so
$h_1=\cL_{X_1}g$ with $(\cL_T+i\sigma)X_1=0$. Put $Y=t_*^{k-1}X_1/(k-1)!$, which is smooth since $t_*$ is regular.
Then $(\cL_T+i\sigma)^{k-1}Y=X_1$, so $h-\cL_Yg$ has order $k-1$.

\emph{Part (2).} Let $(\cL_T)^kh=0$, and split $h=h_0+h_\perp$ into its $\varphi$-average and the rest. The same
induction, using Proposition~\ref{prop:puregauge}(ii), shows that $h_\perp$ is pure gauge. The average $h_0$
satisfies $\cL_{\partial_\varphi}h_0=0$ and $(\cL_{\partial_t})^kh_0=0$. For $k=1$, Proposition~\ref{prop:stationary}
applies. For $k\ge2$:
\begin{itemize}
\item Suppose $\cL_{\partial_t}h_0=c_M\partial_Mg+c_a\partial_ag+\cL_Xg$. The Iyer--Wald charges $\delta H_K$, for
$K\in\{\partial_t,\zeta\}$, are independent of the 2-sphere. They vanish on $\cL_Xg$ for every smooth $X$: from
$\Theta(\cL_Xg)=X\cdot L+dQ_X$ and naturality of $Q$, one gets
$\delta H_K[\cL_Xg]=-\oint(Q_{[X,K]}+Q_{[K,X]})=0$. Applying them on the spheres $\{t_*=\tau\}$ forces
$\delta H_K[c_M\partial_Mg+c_a\partial_ag]=0$. Since $(E,J)=(M/\Xi^2,aM/\Xi^2)$ has Jacobian $M/\Xi^4\ne0$ with
respect to $(M,a)$, it follows that $c_M=c_a=0$. Then $h_0-\cL_{t_*X}g$ has lower order, and the induction
proceeds.
\end{itemize}
\end{proof}

\appendix
\section{Computer-assisted steps}\label{app:cert}\sloppy
All computations use exact arithmetic over $\mathbb Q$ or $\mathbb Q(i)$, except item~C2, which uses rigorous interval
arithmetic. The code is in the repository accompanying this paper; the script names are given in brackets.
\begin{enumerate}[label=C\arabic*.,leftmargin=*]
\item \emph{Dual potential} [\texttt{step39\_c1\_symbolic.py}, \texttt{step89\_c1\_kappa.py}]. After Vieta
elimination, the numerators of $n_1-(\lambda C_0+c_1)$ and $n_2+3$ are divisible by $s^2-4$ as polynomials.
\item \emph{Angular bound} [\texttt{step47\_angular\_proof.py}]. Interval evaluation of the two-level lower bound,
with exact Beta moments, over $(\alpha,\xi)\in[0,0.0718]\times[0,|m|]$ for $|m|\in\{1,2\}$ and $s=\pm2$.
\item \emph{One sign change} [\texttt{step119\_build.py}, \texttt{cert2/p1chart.py}, \texttt{cert2/run\_P1\_all.sh}].
Seventeen polynomials in $(1/t,\zeta)$ of degree at most $11$. Positivity is certified by integer Bernstein-basis
branch and bound on four charts: two rectangles and two blow-up charts at the corner $(1/t,\zeta)=(0,0)$, after
dividing out factors that are positive on the open box. All eight runs certify.
\item \emph{Seam inequality} [\texttt{step121\_build.py}, \texttt{cert2/p2cert.py}, \texttt{cert2/run\_P2\_all.sh}].
Four polynomials in $(1/t,\zeta,S/x_0)$ of degree at most $7$, on nineteen charts covering $[0,1]^3$, including
blow-ups of two degenerate edges. All certify.
\item \emph{Zero frequency} [\texttt{step125\_sigma0\_check.py}]. The Teukolsky--Starobinsky identity at $K=0$ for
the $s=-2$ equation, and the $4\times4$ determinant. For $s=+2$ one substitutes $R_{+2}=\Delta_r^{-2}F$: then
$\Delta_r^{-2}(\Delta_r^3R_{+2}')'=\Delta_r^{-2}(\Delta_rF''-\Delta_r'F'-2\Delta_r''F)$, with
$\Delta_r''=-4\Lambda r^2+2(1-\alpha)$. This matches the $r^2$-terms of the two potentials, $-10\Lambda r^2+8\Lambda r^2=-2\Lambda r^2$,
and the constants are absorbed by the relation between the spin $\pm2$ angular eigenvalues at $\sigma=0$.
\item \emph{Type D identities} [\texttt{step92\_ghp\_check.py}, \texttt{step93\_*.py},\\ \texttt{step95\_pd\_xi.py},
\texttt{step96\_zeta.py}, \texttt{step97\_r5\_pd\_check.py}]. These are checks of identities that are already
derived by hand; the derivations do not depend on them.
\item \emph{Carter ansatz} [\texttt{step102\_carter\_linear.py}]. The Ricci tensor of the ansatz with symbolic
$\Delta_R(R)$ and $\Delta_I(I)$, and the solution of the resulting linear equations.
\item \emph{Schwarzschild--de Sitter Teukolsky--Starobinsky identities}
[\texttt{step126\_sds\_ts.py}, \texttt{step127\_thorn\_ts.py}]. These are the two identities \eqref{eq:sdsTS}. The
first is an operator identity. The second is checked by reducing modulo the radial equation; the remainder
coefficient of $R'$ vanishes identically, and the coefficient of $R$ is the constant $\mathcal C_\ell$. That this
coefficient is independent of $r$ is also a consistency check on the radial potential used.
\end{enumerate}
The Bernstein enclosure routine was tested against random polynomials and random subdivisions
[\texttt{cert2/test\_bern.py}].

\section*{Statement on the use of AI}
In preparing this work I used Claude Opus 5.5 (Anthropic) to develop and check proofs, write and run the symbolic and
certificate code, search and check the literature, and draft and edit the text. The proofs were also checked in
several adversarial review passes, each run by an AI model. These passes found errors, including invalid earlier
versions of the large-$A$ regime of the event lemma and of the passage to the upper half plane, and every error found
was corrected. Numerical experiments were used only to guide the search for proofs, and no result depends on them.
No human expert has yet checked the argument. I have reviewed and edited the content and take full responsibility
for it.

\section*{Data availability}
The code for the computer-assisted steps is available from the author on request.

\begin{thebibliography}{99}
\bibitem{Cha83} S.~Chandrasekhar, \emph{The Mathematical Theory of Black Holes}, Oxford University Press, 1983.
\bibitem{GHP73} R.~Geroch, A.~Held and R.~Penrose, A space-time calculus based on pairs of null directions,
\emph{J. Math. Phys.} \textbf{14} (1973), 874--881.
\bibitem{Hor85} L.~H\"ormander, \emph{The Analysis of Linear Partial Differential Operators III}, Springer, 1985.
\bibitem{Hin26} P.~Hintz, Mode stability for the Klein--Gordon equation on subextremal Kerr--de Sitter spacetimes,
arXiv:2608.25868 (2026).
\bibitem{HPV25} P.~Hintz, O.~Petersen and A.~Vasy, Conditional non-linear stability of Kerr--de Sitter spacetimes in
the full subextremal range, arXiv:2508.06620 (2025).
\bibitem{HV18} P.~Hintz and A.~Vasy, The global non-linear stability of the Kerr--de Sitter family of black holes,
\emph{Acta Math.} \textbf{220} (2018), 1--206.
\bibitem{KS96} A.~Ya.~Kazakov and S.~Yu.~Slavyanov, Integral relations for Heun-class special functions,
\emph{Theoret. Math. Phys.} \textbf{107} (1996), 733--739.
\bibitem{KSFK03} C.~O.~Lousto and B.~F.~Whiting, Reconstruction of black hole metric perturbations from the Weyl
curvature, \emph{Phys. Rev. D} \textbf{66} (2002), 024026.
\bibitem{Kha83} U.~Khanal, Rotating black hole in asymptotic de Sitter space: perturbation of the cosmological
horizon, \emph{Phys. Rev. D} \textbf{28} (1983), 1291--1297.
\bibitem{PV24} O.~Petersen and A.~Vasy, Stationarity and Fredholm theory in subextremal Kerr--de Sitter spacetimes,
\emph{SIGMA} \textbf{20} (2024), 052.
\bibitem{PV25} O.~Petersen and A.~Vasy, Wave equations in the Kerr--de Sitter spacetime: the full subextremal range,
\emph{J. Eur. Math. Soc.} \textbf{27} (2025), 3497--3526.
\bibitem{STU98} H.~Suzuki, E.~Takasugi and H.~Umetsu, Perturbations of Kerr--de Sitter black holes and Heun's
equations, \emph{Prog. Theor. Phys.} \textbf{100} (1998), 491--505.
\bibitem{Tak03} K.~Takemura, Integral representation of solutions to Fuchsian system and Heun's equation,
\emph{J. Math. Anal. Appl.} \textbf{342} (2008), 52--69.
\bibitem{Vas13} A.~Vasy, Microlocal analysis of asymptotically hyperbolic and Kerr--de Sitter spaces,
\emph{Invent. Math.} \textbf{194} (2013), 381--513.
\bibitem{Wal78} R.~M.~Wald, Construction of solutions of gravitational, electromagnetic, or other perturbation
equations from solutions of decoupled equations, \emph{Phys. Rev. Lett.} \textbf{41} (1978), 203--206.
\end{thebibliography}

\end{document}
